\section{AIME I 2009}
        \begin{enumerate}
        	\item (\textit{AIME I 2009}) Call a $3$-digit number \textit{geometric} if it has $3$ distinct digits which, when read from left to right, form a geometric sequence.  Find the difference between the largest and smallest geometric numbers.


 
 \textbf{Answer: }840



\textbf{Solution 1}

Assume that the largest geometric number starts with a $9$. We know that the common ratio must be a rational of the form $k/3$ for some integer $k$, because a whole number should be attained for the 3rd term as well. When $k = 1$, the number is $931$. When $k = 2$, the number is $964$. When $k = 3$, we get $999$, but the integers must be distinct. By the same logic, the smallest geometric number is $124$. The largest geometric number is $964$ and the smallest is $124$. Thus the difference is $964 - 124 = \boxed{840}$.

 





\textbf{Solution 2}

Consider the three-digit number $abc$. If its digits form a geometric progression, we must have that ${a \over b} = {b \over c}$, that is, $b^2 = ac$. 

 The minimum and maximum geometric numbers occur when $a$ is minimized and maximized, respectively. The minimum occurs when $a = 1$; letting $b = 2$ and $c = 4$ achieves this, so the smallest possible geometric number is 124. 

 For the maximum, we have that $b^2 = 9c$; $b$ is maximized when $9c$ is the greatest possible perfect square; this happens when $c = 4$, yielding $b = 6$. Thus, the largest possible geometric number is 964. 

 Our answer is thus $964 - 124 = \boxed{840}$.

 





\textbf{Solution 3}

The smallest geometric number is $124$ because $123$ and any number containing a zero does not work. $964$ is the largest geometric number because the middle digit cannot be 8 or 7. Subtracting the numbers gives $\boxed{840}.$

 




	\item (\textit{AIME I 2009}) There is a complex number $z$ with imaginary part $164$ and a positive integer $n$ such that


 
 \[\frac {z}{z + n} = 4i.\]


 Find $n$.


 
 \textbf{Answer: }697



\textbf{Solution 1}

Let $z = a + 164i$.

 Then
 \[\frac {a + 164i}{a + 164i + n} = 4i\]and
 \[a + 164i = \left (4i \right ) \left (a + n + 164i \right ) = 4i \left (a + n \right ) - 656.\]

 By comparing coefficients, equating the real terms on the leftmost and rightmost side of the equation, 

 we conclude that
 \[a = -656.\]

 By equating the imaginary terms on each side of the equation,

 we conclude that
 \[164i = 4i \left (a + n \right ) = 4i \left (-656 + n \right ).\]

 We now have an equation for $n$: 
 \[4i \left (-656 + n \right ) = 164i,\]

 and this equation shows that $n = \boxed{697}.$

 


 





\textbf{Solution 2}


 \[\frac {z}{z+n}=4i\]

 
 \[1-\frac {n}{z+n}=4i\]

 
 \[1-4i=\frac {n}{z+n}\]

 
 \[\frac {1}{1-4i}=\frac {z+n}{n}\]

 
 \[\frac {1+4i}{17}=\frac {z}{n}+1\]

 Since their imaginary part has to be equal,

 
 \[\frac {4i}{17}=\frac {164i}{n}\]

 
 \[n=\frac {(164)(17)}{4}=697\]

 
 \[n = \boxed{697}.\]

 





\textbf{Solution 3 (Not Highly Recommended)}

Below is an image of the complex plane. Let $\operatorname{Im}(z)$ denote the imaginary part of a complex number $z$.
                    \begin{figure}[H]
                        \centering
                        \includegraphics[width=0.5\textwidth]{Solution/AIME_2009_I_2_Solution_Figure_1.png}
                    \end{figure}
                    $z$ must lie on the line $\operatorname{Im}(z)=164$. $z+n$ must also lie on the same line, since $n$ is real and does not affect the imaginary part of $z$.

 Consider $z$ and $z+n$ in terms of their magnitude (distance from the origin) and phase (angle formed by the point, the origin, and the positive real axis, measured counterclockwise from said axis). When multiplying/dividing two complex numbers, you can multiply/divide their magnitudes and add/subtract their phases to get the magnitude and phase of the product/quotient. Expressed in a formula, we have $z_1z_2 = r_1\angle\theta_1 \cdot r_2\angle\theta_2 = r_1r_2\angle(\theta_1+\theta_2)$ and $\frac{z_1}{z_2} = \frac{r_1\angle\theta_1}{r_2\angle\theta_2} = \frac{r_1}{r_2}\angle(\theta_1-\theta_2)$, where $r$ is the magnitude and $\theta$ is the phase, and $z_n=r_n\angle\theta_n$.

 Since $4i$ has magnitude $4$ and phase $90^\circ$ (since the positive imaginary axis points in a direction $90^\circ$ counterclockwise from the positive real axis), $z$ must have a magnitude $4$ times that of $z+n$. We denote the length from the origin to $z+n$ with the value $x$ and the length from the origin to $z$ with the value $4x$. Additionally, $z$, the origin, and $z+n$ must form a right angle, with $z$ counterclockwise from $z+n$.

 This means that $z$, the origin, and $z+n$ form a right triangle. The hypotenuse is the length from $z$ to $z+n$ and has length $n$, since $n$ is defined to be a positive integer. The area of the triangle can be expressed using the two legs, as $\frac{x \cdot 4x}{2}$, or using the hypotenuse and its corresponding altitude, as $\frac{164n}{2}$, so $\frac{x \cdot 4x}{2} = \frac{164n}{2} \implies x^2 = 41n$. By Pythagorean Theorem, $x^2+(4x)^2 = n^2 \implies 17x^2 = n^2$. Substituting out $x^2$ using the earlier equation, we get $17\cdot41n = n^2 \implies n = \boxed{697}$. ~\href{https://artofproblemsolving.com/wiki/index.php?title=User:Emerald\_block}{emerald\_block}

 





\textbf{Solution 4 (fast)}

Taking the reciprocal of our equation gives us $1 + \frac{n}{z} = \frac{1}{4i}.$ Therefore, 
 \[\frac{n}{z} = \frac{1-4i}{4i} = \frac{17}{-16+4i}.\] Since $z$ has an imaginary part of $164$, we must multiply both sides of our RHS fraction by $\frac{164}{4} = 41$ so that its denominator's imaginary part matches the LHS fraction's denominator's imaginary part. That looks like this:
 \[\frac{n}{z} = \frac{697}{-656 + 164i}.\] Therefore, we can conclude the real part of $z$ is $-656$ and $n = \boxed{697}.$ (it wasn't necessary to find the real part)

 




	\item (\textit{AIME I 2009}) A coin that comes up heads with probability $p > 0$ and tails with probability $1 - p > 0$ independently on each flip is flipped eight times. Suppose the probability of three heads and five tails is equal to $\frac {1}{25}$ of the probability of five heads and three tails. Let $p = \frac {m}{n}$, where $m$ and $n$ are relatively prime positive integers. Find $m + n$.


 
 \textbf{Answer: }011



\textbf{Solution 1}

The probability of three heads and five tails is $\binom {8}{3}p^3(1-p)^5$ and the probability of five heads and three tails is $\binom {8}{3}p^5(1-p)^3$.

 
 \begin{align*} 25\binom {8}{3}p^3(1-p)^5&=\binom {8}{3}p^5(1-p)^3 \\ 25(1-p)^2&=p^2 \\ 25p^2-50p+25&=p^2 \\ 24p^2-50p+25&=0 \\ p&=\frac {5}{6}\end{align*}


 Therefore, the answer is $5+6=\boxed{011}$.

 





\textbf{Solution 2}

We start as shown above. However, when we get to $25(1-p)^2=p^2$, we square root both sides to get $5(1-p)=p$. We can do this because we know that both $p$ and $1-p$ are between $0$ and $1$, so they are both positive. Now, we have: 

 
 \begin{align*} 5(1-p)&=p \\ 5-5p&=p \\ 5&=6p \\ p&=\frac {5}{6}\end{align*}


 Now, we get $5+6=\boxed{011}$.

 





\textbf{Solution 3}

Rewrite it as : $(P)^3$$(1-P)^5=\frac {1}{25}$ $(P)^5$$(1-P)^3$

 This can be simplified as $24P^2 -50P + 25 = 0$

 This can be factored into $(4P-5)(6P-5)$

 This yields two solutions: $\frac54$ (ignored because it would result in $1-p<0$  ) or $\frac56$

 Therefore, the answer is $5+6$ = $\boxed {011}$

 




	\item (\textit{AIME I 2009}) In parallelogram $ABCD$, point $M$ is on $\overline{AB}$ so that $\frac {AM}{AB} = \frac {17}{1000}$ and point $N$ is on $\overline{AD}$ so that $\frac {AN}{AD} = \frac {17}{2009}$. Let $P$ be the point of intersection of $\overline{AC}$ and $\overline{MN}$. Find $\frac {AC}{AP}$.


 
 \textbf{Answer: }177



\textbf{Solution 1 (CHEESE!)}

One of the ways to solve this problem is to make this parallelogram a straight line.So the whole length of the line is $APC$($AMC$ or $ANC$), and $ABC$ is $1000x+2009x=3009x.$

 $AP$($AM$ or $AN$) is $17x.$

 So the answer is $3009x/17x = \boxed{177}$

 





\textbf{Solution 2}

Draw a diagram with all the given points and lines involved. Construct parallel lines $\overline{DF_2F_1}$ and $\overline{BB_1B_2}$ to $\overline{MN}$, where for the lines the endpoints are on $\overline{AM}$ and $\overline{AN}$, respectively, and each point refers to an intersection. Also, draw the median of quadrilateral $BB_2DF_1$ $\overline{E_1E_2E_3}$ where the points are in order from top to bottom. Clearly, by similar triangles, $BB_2 = \frac {1000}{17}MN$ and $DF_1 = \frac {2009}{17}MN$. It is not difficult to see that $E_2$ is the center of quadrilateral $ABCD$ and thus the midpoint of $\overline{AC}$ as well as the midpoint of $\overline{B_1}{F_2}$ (all of this is easily proven with symmetry). From more triangle similarity, $E_1E_3 = \frac12\cdot\frac {3009}{17}MN\implies AE_2 = \frac12\cdot\frac {3009}{17}AP\implies AC = 2\cdot\frac12\cdot\frac {3009}{17}AP$$= \boxed{177}AP$.

 





\textbf{Solution 3}

Using vectors, note that $\overrightarrow{AM}=\frac{17}{1000}\overrightarrow{AB}$ and $\overrightarrow{AN}=\frac{17}{2009}\overrightarrow{AD}$.  Note that $\overrightarrow{AP}=\frac{x\overrightarrow{AM}+y\overrightarrow{AN}}{x+y}$ for some positive x and y, but at the same time is a scalar multiple of $\overrightarrow{AB}+\overrightarrow{AD}$.  So, writing the equation $\overrightarrow{AP}=\frac{x\overrightarrow{AM}+y\overrightarrow{AN}}{x+y}$ in terms of $\overrightarrow{AB}$ and $\overrightarrow{AD}$, we have $\overrightarrow{AP}=\frac{\frac{17x}{1000}\overrightarrow{AB}+\frac{17y}{2009}\overrightarrow{AD}}{x+y}$.  But the coefficients of the two vectors must be equal because, as already stated, $\overrightarrow{AP}$ is a scalar multiple of $\overrightarrow{AB}+\overrightarrow{AD}$.  We then see that $\frac{x}{x+y}=\frac{1000}{3009}$ and $\frac{y}{x+y}=\frac{2009}{3009}$. Finally, we have $\overrightarrow{AP}=\frac{17}{3009}(\overrightarrow{AB}+\overrightarrow{AD})$ and, simplifying, $\overrightarrow{AB}+\overrightarrow{AD}=177\overrightarrow{AP}$ and the desired quantity is $177$.

 





\textbf{Solution 4}

We approach the problem using mass points on triangle $ABD$ as displayed below.
                    \begin{figure}[H]
                        \centering
                        \includegraphics[width=0.5\textwidth]{Solution/AIME_2009_I_4_Solution_Figure_1.png}
                    \end{figure}
                    

 But as $MN$ does not protrude from a vertex, we will have to "split the mass" at point $A$. First, we know that $DO$  is congruent to $BO$ because diagonals of parallelograms bisect each other. Therefore, we can assign equal masses to points $B$ and $D$. In this case, we assign $B$ and $D$ a mass of 17 each. Now we split the mass at $A$, so we balance segments $AB$ and $AD$ separately, and then the mass of $A$ is the sum of those masses. A mass of 983 is required to balance segment $AB$, while a mass of 1992 is required to balance segment $AD$. Therefore, $A$ has a mass of $1992+983=2975$. Also, $O$ has a mass of 34.  Therefore, $\frac{AO}{AP}=\frac{2975+34}{34}=\frac{3009}{34}$, so $\frac{AC}{AP}=\frac{2 (3009)}{34}=177$.

 





\textbf{Solution 5}

Assume, for the ease of computation, that $AM=AN=17$, $AB=1000$, and $AD=2009$.  Now, let line $MN$ intersect line $CD$ at point $X$ and let $Y$ be a point such that $XY\parallel AD$ and $AY\parallel DX$.  As a result, $ADXY$ is a parallelogram.  By construction, $\triangle MAN\sim \triangle MYX$ so 
 \[\frac{MY}{MA}=\frac{YX}{AN}=\frac{AD}{AN}=\frac{2009}{17}\implies MY=2009\] and $AY=DX=2009-17$.  Also, because $AM\parallel XC$, we have $\triangle PAM\sim \triangle PCX$ so 
 \[\frac{PC}{PA}=\frac{CX}{AM}=\frac{DX+CD}{AM}=\frac{2009-17+1000}{17}=176.\]  Hence, $\frac{AC}{AP}=\frac{PC}{PA}+1=\boxed{177}.$

 





\textbf{Solution 6(Coordinate Geometry)}

Assign $A = (0,0)$. Since there are no constraints in the problem against this, assume $ABCD$ to be a rectangle with dimensions $1000 \times 2009.$ Now, we can assign 
 \[A=(0,0)\] 
 \[B=(1000, 0)\] 
 \[C=(1000,-2009)\] 
 \[D=(0,-2009).\]

 Then, since $\frac{AM}{AB} = \frac{17}{1000}$ and $AB = 1000$, we can place $M$ at $(17, 0).$ Similarly, place $AN$ at $(0, 17).$Then, the equation of line $MN$ is $y=x-17,$ and the equation of $AC$ is $y=\frac{-2009}{1000}x.$ Solve to find point $P$ at $\left( \frac{1000}{177}, \frac{-2009}{177} \right)$.

 We can calculate vectors to represent the distances: 
 \[\overrightarrow{AC}= <1000, -2009>\] 
 \[\overrightarrow{AP}= \frac{1}{177}<1000, -2009>.\] In this way, we can see that 
 \[AC:AP = 177:1,\] and our answer is $\boxed{177}.$

 \textbf{Note} It is possible to use \href{https://artofproblemsolving.com/wiki/index.php?title=Coordinate\_geometry}{coordinate geometry} without setting $ABCD$ as a rectangle, either by projecting the plane onto another (tilted) plane or removing the restriction that the \href{https://artofproblemsolving.com/wiki/index.php?title=Axes}{axes} have to be \href{https://artofproblemsolving.com/wiki/index.php?title=Perpendicular}{perpendicular}.

 





\textbf{Solution 7 (Coordinates and Similar triangles)}

Let $ABCD$ be a square with side length $1$ where $D = (0,0)$, $A = (0,1)$, $B = (1,1)$, and $C = (1,0)$.

 Draw $M$ on $AB$ such that $AM$ has length $\frac{17}{1000}$ and $AN$ has length $\frac{17}{2009}$. Draw $AC$ with $P$ as the intersection point of $AC$ and $MN$.

 $N$ has coordinates $(0, 1-\frac{17}{2009}) = (0, \frac{1992}{2009})$

 Extend lines $MN$ and $CD$ such that their intersection point is point $E$, which lies on $y=0$.

 Line $ME$ has slope $\frac{AN}{AM} = \frac{1000}{2009}$. 

 With the y-intercept $N$ it has the equation $y = \frac{1000}{2009}x + \frac{1992}{2009}$.

 Solving for the $x$ coordinate on point $E$ $(y = 0)$, $x = -\frac{1992}{1000}$. $ED$ has length $\frac{1992}{1000}$ and $EC$ has length $\frac{1992}{1000} + 1 = \frac{2992}{1000}$. 

 Triangles $APM$ and $CPE$ are similar (AA). $\frac{AP}{PC} = \frac{AM}{EC} = \frac{2992}{17}$. $PC = \frac{2992}{17}AP$

 $\frac{AC}{AP} = \frac{PC+AP}{AP} = \frac{3009}{17} = 177$.

 




	\item (\textit{AIME I 2009}) Triangle $ABC$ has $AC = 450$ and $BC = 300$. Points $K$ and $L$ are located on $\overline{AC}$ and $\overline{AB}$ respectively so that $AK = CK$, and $\overline{CL}$ is the angle bisector of angle $C$. Let $P$ be the point of intersection of $\overline{BK}$ and $\overline{CL}$, and let $M$ be the point on line $BK$ for which $K$ is the midpoint of $\overline{PM}$. If $AM = 180$, find $LP$.


 
 \textbf{Answer: }072



\textbf{Solution 1}


                    \begin{figure}[H]
                        \centering
                        \includegraphics[width=0.5\textwidth]{Solution/AIME_2009_I_5_Solution_Figure_1.png}
                    \end{figure}
                    Since $K$ is the midpoint of $\overline{PM}$ and $\overline{AC}$, quadrilateral $AMCP$ is a parallelogram, which implies $AM||LP$ and  $\bigtriangleup{AMB}$ is similar to $\bigtriangleup{LPB}$

 Thus,

 
 \[\frac {AM}{LP}=\frac {AB}{LB}=\frac {AL+LB}{LB}=\frac {AL}{LB}+1\]

 Now let's apply the angle bisector theorem.

 
 \[\frac {AL}{LB}=\frac {AC}{BC}=\frac {450}{300}=\frac {3}{2}\]

 
 \[\frac {AM}{LP}=\frac {AL}{LB}+1=\frac {5}{2}\]

 
 \[\frac {180}{LP}=\frac {5}{2}\]

 
 \[LP=\boxed {072}\]

 \begin{verbatim}
Note: Another way to realize that $\bigtriangleup{AMB}$ is similar to $\bigtriangleup{LPB}$ is through angle chasing and congruent triangles. Angle $AKM$ is congruent to angle $CKP$. Segments $MK$ and $PK$ are congruent as stated in the problem. Segments $AK$ and $CK$ are congruent as stated in the problem. Thus, by Side-Angle-Side congruency $\bigtriangleup{AMK}$ is similar to $\bigtriangleup{CPK}$. Then angle $AMK = AMB$ is congruent to angle $CPK = BPL$. Notice triangles $\bigtriangleup{AMB}$ and $\bigtriangleup{LPB}$ both share angle $MBA$. Thus the two triangles are similar because of AA similarity with angles [$AMB$ and $LPB$] and [$MBA = PBL$]. *Solution in Beauty of Math video written out
\end{verbatim}





\textbf{Solution 2 (Mass Points)}

Using the diagram above, we can solve this problem by using mass points. By angle bisector theorem:
 \[\frac{BL}{CB}=\frac{AL}{CA}\implies\frac{BL}{300}=\frac{AL}{450}\implies 3BL=2AL\]So, we can weight $A$ as $2$ and $B$ as $3$ and $L$ as $5$. Since $K$ is the midpoint of $A$ and $C$, the weight of $A$ is equal to the weight of $C$, which equals $2$.Also, since the weight of $L$ is $5$ and $C$ is $2$, we can weight $P$ as $7$.

 By the definition of mass points, 
 \[\frac{LP}{CP}=\frac{2}{5}\implies LP=\frac{2}{5}CP\]By vertical angles, angle $MKA =$ angle $PKC$. Also, it is given that $AK=CK$ and $PK=MK$.

 By the SAS congruence, $\triangle MKA$ = $\triangle PKC$. So, $MA$ = $CP$ = $180$.Since $LP=\frac{2}{5}CP$, $LP = \frac{2}{5}(180) = \boxed{072}$

 





\textbf{Solution 3 (Law of Cosines Bash)}

Using the diagram from solution $1$, we can also utilize the fact that $AMCP$ forms a parallelogram. Because of that, we know that $AM = CP = 180$.

 Applying the angle bisector theorem to $\triangle CKB$, we get that $\frac{KP}{PB} = \frac{225}{300} = \frac{3}{4}.$ So, we can let $MK = KP = 3x$ and $BP = 4x$. 

 Now, apply law of cosines on $\triangle CKP$ and $\triangle CPB.$

 If we let $\angle KCP = \angle PCB = \alpha$, then the law of cosines gives the following system of equations: 

 
 \[9x^2 = 225^2 + 180^2 - 2\cdot 225 \cdot 180 \cdot \cos \alpha\]
 \[16x^2 = 180^2 + 300^2 - 2 \cdot 180 \cdot 300 \cdot \cos \alpha.\]

 Bashing those out, we get that $x = 15 \sqrt{13}$ and $\cos \alpha = \frac{7}{10}.$

 Since $\cos \alpha = \frac{7}{10}$, we can use the double angle formula to calculate that $\cos \left(2 \alpha \right) = -\frac{1}{50}.$

 Now, apply Law of Cosines on $\triangle ABC$ to find $AB$. 

 We get: 
 \[AB^2 = 450^2 + 300^2 - 2 \cdot 450 \cdot 300 \cdot \left(- \frac{1}{50} \right).\]

 Bashing gives $AB = 30 \sqrt{331}.$

 From the angle bisector theorem on $\triangle ABC$, we know that $\frac{AL}{BL} = \frac{450}{300} = \frac{3}{2}.$ So, $AL = 18 \sqrt{331}$ and $BL = 12 \sqrt{331}.$

 Now, we apply Law of Cosines on $\triangle ALC$ and $\triangle BLC$ in order to solve for the length of $LC$. 

 We get the following system: 

 
 \[(18 \sqrt{331})^2 = 450^2 + LC^2 - 2 \cdot 450 \cdot LC \cdot \frac{7}{10}\]
 \[(12 \sqrt{331})^2 = LC^2 + 300^2 - 2 \cdot 300 \cdot LC \cdot \frac{7}{10}\]

 The first equation gives $LC = 252$ or $378$ and the second gives $LC = 252$ or $168$.

 The only value that satisfies both equations is $LC = 252$, and since $LP = LC - PC$, we have 
 \[LC = 252 - 180 = \boxed{072}.\]

 





\textbf{Solution 4(Area Ratios)}

Note that we are given that $\overline{MK} = \overline{KP}$, that $\overline{AK} = \overline{CK}.$ Note then that $\angle MKA = \angle CKB$ by vertical angles. From this, we have $\triangle MKA \cong PKC.$ This means that $\overline{CP}$ is 180. Applying angle bisector theorem on $\triangle ACB$ gives $\frac{\overline{AL}}{\overline{LB}} = \frac{450}{300} = \frac{3}{2}.$ Applying it on $\triangle KCB$ yields
 \[\frac{\overline{KP}}{\overline{PB}} = \frac{225}{300} = \frac{3}{4}\]Now we can proceed with area ratios. Suppose the area of $\triangle ACB = A.$ This means that 
 \[[\triangle AKL] = \left(\frac{225}{225+225}\right)\left(\frac{3}{5}\right)A = \frac{3}{10}A\]Continuing on $\triangle LPB$ we have 
 \[[\triangle LPB] = \left(\frac{2}{2+3}\right)\left(\frac{1}{2}\right)\left(\frac{4}{4+3}\right) = \frac{4}{35}A\]Since $\overline{AK}=\overline{KC}$ $[\triangle KPL] = [\triangle AKB]-[\triangle AKL] - [\triangle LPB] = \frac{1}{2}A - \frac{3}{10}A - \frac{4}{35}A = \frac{3}{35}A.$Area ratios on $\triangle KCP$ yield $[\triangle KCP] = \left(\frac{1}{2}\right)\left(\frac{3}{3+4}\right) = \frac{3}{14}.$ Now, suppose $\overline{LP} = x.$ We have that the ratio of areas of $\triangle LKP$ and $\triangle PKC$ is $\frac{x}{180}$and is also $\frac{\frac{3}{35}}{\frac{3}{14}}$ and equating these gives 
 \[x = \boxed{72}\]

 




	\item (\textit{AIME I 2009}) How many positive integers $N$ less than $1000$ are there such that the equation $x^{\lfloor x\rfloor} = N$ has a solution for $x$?


 
 \textbf{Answer: }412



\textbf{Solution 1}

First, $x$ must be less than $5$, since otherwise $x^{\lfloor x\rfloor}$ would be at least $3125$ which is greater than $1000$. 

 Because ${\lfloor x\rfloor}$ must be an integer, let’s do case work based on ${\lfloor x\rfloor}$:

 For ${\lfloor x\rfloor}=0$, $N=1$ as long as $x \neq 0$. This gives us $1$ value of $N$.

 For ${\lfloor x\rfloor}=1$, $N$ can be anything between $1^1$ to $2^1$ excluding $2^1$

 Therefore, $N=1$. However, we got $N=1$ in case 1 so it got counted twice.

 For ${\lfloor x\rfloor}=2$, $N$ can be anything between $2^2$ to $3^2$ excluding $3^2$

 This gives us $3^2-2^2=5$ $N$'s

 For ${\lfloor x\rfloor}=3$, $N$ can be anything between $3^3$ to $4^3$ excluding $4^3$

 This gives us $4^3-3^3=37$ $N$'s

 For ${\lfloor x\rfloor}=4$, $N$ can be anything between $4^4$ to $5^4$ excluding $5^4$

 This gives us $5^4-4^4=369$ $N$'s

 Since $x$ must be less than $5$, we can stop here and the answer is $1+5+37+369= \boxed {412}$ possible values for $N$.

 Alternatively, one could find that the values which work are $1^1,\ 2^2,\ 3^3,\ 4^4,\ \sqrt{5}^{\lfloor\sqrt{5}\rfloor},\ \sqrt{6}^{\lfloor\sqrt{6}\rfloor},\ \sqrt{7}^{\lfloor\sqrt{7}\rfloor},\ \sqrt{8}^{\lfloor\sqrt{8}\rfloor},\ \sqrt[3]{28}^{\lfloor\sqrt[3]{28}\rfloor},\ \sqrt[3]{29}^{\lfloor\sqrt[3]{29}\rfloor},\ \sqrt[3]{30}^{\lfloor\sqrt[3]{30}\rfloor},\ ...,$$\ \sqrt[3]{63}^{\lfloor\sqrt[3]{63}\rfloor},\ \sqrt[4]{257}^{\lfloor\sqrt[4]{257}\rfloor},\ \sqrt[4]{258}^{\lfloor\sqrt[4]{258}\rfloor},\ ...,\ \sqrt[4]{624}^{\lfloor\sqrt[4]{624}\rfloor}$ to get the same answer.

 





\textbf{Solution 2}

For a positive integer $k$, we find the number of positive integers $N$ such that $x^{\lfloor x\rfloor}=N$ has a solution with ${\lfloor x\rfloor}=k$. Then $x=\sqrt[k]{N}$, and because $k \le x < k+1$, we have $k^k \le x^k < (k+1)^k$, and because $(k+1)^k$ is an integer, we get $k^k \le x^k \le (k+1)^k-1$. The number of possible values of $x^k$ is equal to the number of integers between $k^k$ and $(k+1)^k-1$ inclusive, which is equal to the larger number minus the smaller number plus one or $((k+1)^k-1)-(k^k)+1$, and this is equal to $(k+1)^k-k^k$. If $k>4$, the value of $x^k$ exceeds $1000$, so we only need to consider $k \le 4$. The requested number of values of $N$ is the same as the number of values of $x^k$, which is $\sum^{4}_{k=1} [(k+1)^k-k^k]=2-1+9-4+64-27+625-256=\boxed{412}$.

 \textbf{Video Solution 1} \textbf{Video Solution 2} \textbf{Video Solution 3} 




	\item (\textit{AIME I 2009}) The sequence $(a_n)$ satisfies $a_1 = 1$ and $5^{(a_{n + 1} - a_n)} - 1 = \frac {1}{n + \frac {2}{3}}$ for $n \geq 1$. Let $k$ be the least integer greater than $1$ for which $a_k$ is an integer. Find $k$.


 
 \textbf{Answer: }041



\textbf{Solution}

The best way to solve this problem is to get the iterated part out of the exponent:
 \[5^{(a_{n + 1} - a_n)} = \frac {1}{n + \frac {2}{3}} + 1\]
 \[5^{(a_{n + 1} - a_n)} = \frac {n + \frac {5}{3}}{n + \frac {2}{3}}\]
 \[5^{(a_{n + 1} - a_n)} = \frac {3n + 5}{3n + 2}\]
 \[a_{n + 1} - a_n = \log_5{\left(\frac {3n + 5}{3n + 2}\right)}\]
 \[a_{n + 1} - a_n = \log_5{(3n + 5)} - \log_5{(3n + 2)}\]Plug in $n = 1, 2, 3, 4$ to see the first few terms of the sequence: 
 \[\log_5{5},\log_5{8}, \log_5{11}, \log_5{14}.\] We notice that the terms  $5, 8, 11, 14$ are in arithmetic progression. Since $a_1 = 1 = \log_5{5} = \log_5{(3(1) + 2)}$, we can easily use induction to show that $a_n = \log_5{(3n + 2)}$. So now we only need to find the next value of $n$ that makes $\log_5{(3n + 2)}$ an integer. This means that $3n + 2$ must be a power of $5$. We test $25$: 
 \[3n + 2 = 25\] 
 \[3n = 23\] This has no integral solutions, so we try $125$: 
 \[3n + 2 = 125\] 
 \[3n = 123\] 
 \[n = \boxed{041}\]

 





\textbf{Solution 2 (Telescoping)}

We notice that by multiplying the equation from an arbitrary $a_n$ all the way to $a_1$, we get:
 \[5^{a_n-a_1}=\dfrac{n+\frac23}{1+\frac23}\]This simplifies to
 \[5^{a_n}=3n+2.\]We can now test powers of $5$.

 $5$ - that gives us $n=1$, which is useless.

 $25$ - that gives a non-integer $n$.

 $125$ - that gives $n=\boxed{41}$.

 -integralarefun

 





\textbf{Solution 3 (I did too many FE's)}

The given recursion is equivalent to $\frac{5^{a_{n+1}}}{3(n+1)+2}=\frac{5^{a_n}}{3n+2}$. By a simple inductive argument, the value $\frac{5^{a_k}}{3k+2}$ is constant for all naturals $k$. In particular, when $k=1$, the expression is equal to $1$, so $\frac{5^{a_k}}{3k+2}=1$ for all naturals $k$. Therefore, $a_k=\log_5(3k+2)$, and testing powers of $5$ yields $k=\boxed{41}$ as the answer, which is when $a_{41}=\log_5(3\cdot41+2)=\log_5(125)=3$.

 





\textbf{Solution 4 (Easier Explicit)}

Note that 

 
 \[5^{a_{n+1} - a_n} = \frac{3n+5}{3n+2}\]
 \[a_{n+1} - a_n = \log_{5}(\frac{3n+5}{3n+2})\]
 \[a_{n+1} = \log_{5}(\frac{3n+5}{3n+2}) + a_n\].

 So,

 
 \[a_{n+1} = \log_{5}(\frac{3n+5}{3n+2}) + \log_{5}(\frac{3(n-1)+5}{3(n-1)+2}) + \dots +  \log_{5}(\frac{8}{5}) + a_1\]

 in which we see telescopes to 

 
 \[a_{n+1} = \log_{5}(\frac{3n+5}{5})\]

 Now, we must find the smallest $n$ for which $\log_{5}(\frac{3n+5}{5}) \in \mathbb{Z} \neq 1, n \in \mathbb{Z}^+ \neq 0$. We first check $\frac{3n+5}{5} = 5$, but see that $n$ is not an integer, so we check $\frac{3n+5}{5} = 25$, so $3n = 120$, $n = 40$. The index $k = n+1$ implies $k = \boxed{041}$.

 




	\item (\textit{AIME I 2009}) Let $S = \{2^0,2^1,2^2,\ldots,2^{10}\}$. Consider all possible positive differences of pairs of elements of $S$. Let $N$ be the sum of all of these differences. Find the remainder when $N$ is divided by $1000$.


 
 \textbf{Answer: }398



\textbf{Solution 1 (bash)}

When computing $N$, the number $2^x$ will be added $x$ times (for terms $2^x-2^0$, $2^x-2^1$, ..., $2^x - 2^{x-1}$), and subtracted $10-x$ times. Hence $N$ can be computed as $N=10\cdot 2^{10} + 8\cdot 2^9 + 6\cdot 2^8 + \cdots - 8\cdot 2^1 - 10\cdot 2^0$. Evaluating $N \bmod {1000}$ yields:

 
 \begin{align*} N  & = 10(2^{10}-1) + 8(2^9 - 2^1) + 6(2^8-2^2) + 4(2^7-2^3) + 2(2^6-2^4) \\ & = 10(1023) + 8(510) + 6(252) + 4(120) + 2(48) \\ & = 10(1000+23) + 8(500+10) + 6(250+2) + 480 + 96 \\ & \equiv (0 + 230) + (0 + 80) + (500 + 12) + 480 + 96 \\ & \equiv \boxed{398} \end{align*}


 





\textbf{Solution 2}

This solution can be generalized to apply when $10$ is replaced by other positive integers.

 Extending from Solution 1, we get that the sum $N$ of all possible differences of pairs of elements in $S$ when $S = \{2^0,2^1,2^2,\ldots,2^{n}\}$ is equal to $\sum_{x=0}^{n} (2x-n) 2^x$. Let $A = \sum_{x=0}^{n} x2^x$, $B=\sum_{x=0}^{n} 2^x$. Then $N=2A - nB$.

 For $n = 10$, $B = 2^{11}-1 = 2047 \equiv 47 \pmod{1000}$ by the geometric sequence formula. 

 $2A = \sum_{x=1}^{n+1} (x-1)2^x$, so $2A - A = A = n2^{n+1} - \sum_{x=1}^{n} 2^x$. Hence, for $n = 10$, $A = 10 \cdot 2^{11} - 2^{11} + 2 = 9 \cdot 2^{11} + 2 \equiv 48 \cdot 9 + 2 =$

 $434 \pmod{1000}$, by the geometric sequence formula and the fact that $2^{10} = 1024 \equiv 24 \pmod{1000}$.

 Thus, for $n = 10$, $N = 2A - 10B \equiv 2\cdot 434 - 10\cdot 47 = 868 - 470 = \boxed{398}$.

 





\textbf{Solution 3}

Consider the unique differences $2^{a + n} - 2^a$. Simple casework yields a sum of $\sum_{n = 1}^{10}(2^n - 1)(2^{11 - n} - 1) = \sum_{n = 1}^{10}2^{11} + 1 - 2^n - 2^{11 - n} = 10\cdot2^{11} + 10 - 2(2 + 2^2 + \cdots + 2^{10})$$= 10\cdot2^{11} + 10 - 2^2(2^{10} - 1)\equiv480 + 10 - 4\cdot23\equiv\boxed{398}\pmod{1000}$. This method generalizes nicely as well.

 





\textbf{Solution 4 (Extreme bash)}

Find the positive differences in all $55$ pairs and you will get $\boxed{398}$. (This is not recommended unless you can't find any other solutions to this problem)

 \textit{List of Differences}:

 1, 3, 7, 15, 31, 63, 127, 255, 511, 1023, 2, 6, 14, 30, 62, 126, 254, 510, 1022, 4, 12, 28, 60, 124, 252, 508, 1020, 8, 24, 56, 120, 248, 504, 1016, 16, 48, 112, 240, 496, 1008, 32, 96, 224, 480, 992, 64, 192, 448, 960, 128, 384, 896, 256, 768, 512

 





\textbf{Solution 5 (Quite Less Bash)}

Every value $2^a - 2^b$ is guaranteed to be different if for all pairs of numbers $(a,b)$ is unordered. In other terms, every pair only appears once, with $(a,b)$ being the same as $(b,a)$.

 Imagine cyclically pairing every single element of this list together. Then, we appear with $10$ copies of $2^{10}$, and then from the remaining $9$ differences consisting of $2^9$, when summed, coincide with one difference from $2^{10}$, giving $10-2 = 8$ copies of $2^9$.

 If the above wording is weird to you, think of it like this. One of the pairs consisting of $2^{10} - 2^b$ for some $b$ is $2^{10} - 2^9$. Then, in all the pairs consisting of $2^9 - 2^z$ for some $z$, we see $z \neq 10$ as that absolute value yields the same sum. Therefore, we have the number of copies of $2^{10}$, subtracted with the case including $2^9$, and that result is $10-1 = 9$ copies of $2^9$. However, when we add the cases together, one pair from $2^{10}$ coincides with one pair from $2^9$, specifically, any $2^9$ case and the case $2^{10} - 2^9$, in which one more pair cancels out, giving us 8 copies in the sum. 

 One can then, either through \href{https://artofproblemsolving.com/wiki/index.php?title=Engineer\%27s\_Induction&action=edit&redlink=1}{Engineer's Induction} or through webbing (the process of mapping every case to one another and showing a result must be true), can realize that every preceding case decreases by 2. Then, we obtain

 
 \[\sum_{k=0}^{10} (10-2k)(2^{10-k})\]
 \[10(2^{10}) + 8(2^9) + 6(2^8) + 4(2^7) + 2(2^6) + 0 - 2(2^4) - 4(2^3) - 6(2^2) - 8(2^1) - 10 \pmod{1000}\]
 \[240 + 96 + 536 + 512 + 128 - 32 - 32 - 24 - 16 - 10 \pmod{1000}\]
 \[240 + 96 + 48 + 128 - 32 - 32 - 24 - 16 - 10 \pmod{1000}\]
 \[512 - 114 \pmod{1000}\]
 \[\boxed{398} \pmod{1000}\]

 




	\item (\textit{AIME I 2009}) A game show offers a contestant three prizes A, B and C, each of which is worth a whole number of dollars from $\text{\textdollar}1$ to $\text{\textdollar}9999$ inclusive. The contestant wins the prizes by correctly guessing the price of each prize in the order A, B, C. As a hint, the digits of the three prices are given. On a particular day, the digits given were $1, 1, 1, 1, 3, 3, 3$. Find the total number of possible guesses for all three prizes consistent with the hint.


 
 \textbf{Answer: }420



\textbf{Solution 1}

[Clarification: You are supposed to find the number of all possible tuples of prices, $(A, B, C)$, that could have been on that day.]

 Since we have three numbers, consider the number of ways we can put these three numbers together in a string of 7 digits. For example, if $A=113, B=13, C=31$, then the string is

 
 \[1131331.\]

 Since the strings have seven digits and three threes, there are $\binom{7}{3}=35$ arrangements of all such strings.

 In order to obtain all combination of A,B,C, we partition all the possible strings into 3 groups.

 Let's look at the example. We have to partition it into 3 groups with each group having at least 1 digit. In other words, we need to find the solution to

 
 \[x+y+z=7, x,y,z>0.\]

 This gives us

 
 \[\binom{6}{2}=15\]

 ways by stars and bars. But we have counted the one with 5 digit numbers; that is, $(5,1,1),(1,1,5),(1,5,1)$.

 Thus, each arrangement has 
 \[\binom{6}{2}-3=12\] ways per arrangement, and there are $12\times35=\boxed{420}$ ways.

 





\textbf{Solution 1a (Casework)}

Follow Solution 1 until the partitioning of all the possible strings into 3 groups. Another way to partition the strings is by casework. 

 Case 1: one of [A, B, or C] has one digit, another has two digits, and the last has four digits. There are $3!=6$ ways this can happen.

 Case 2: one of [A, B, or C] has two digit, another has two digits, and the last has three digits.There are $3!/2=3$ ways this can happen.

 Case 3: one of [A, B, or C] has one digit, another has three digits, and the last has three digits.There are $3!/2=3$ ways this can happen.

 The total numbers of ways per arrangement is $6+3+3=12$ ways. Following Solution 1, there are $12\times35=\boxed{420}$ total ways.

 -unhappyfarmer

 




	\item (\textit{AIME I 2009}) The Annual Interplanetary Mathematics Examination (AIME) is written by a committee of five Martians, five Venusians, and five Earthlings.  At meetings, committee members sit at a round table with chairs numbered from $1$ to $15$ in clockwise order.  Committee rules state that a Martian must occupy chair $1$ and an Earthling must occupy chair $15$.  Furthermore, no Earthling can sit immediately to the left of a Martian, no Martian can sit immediately to the left of a Venusian, and no Venusian can sit immediately to the left of an Earthling.  The number of possible seating arrangements for the committee is $N \cdot (5!)^3$.  Find $N$.


 
 \textbf{Answer: }346



\textbf{Solution 1}

Since the 5 members of each planet committee are distinct we get that the number of arrangement of sittings is in the form $N*(5!)^3$ because for each $M, V, E$ sequence we have $5!$ arrangements within the Ms, Vs, and Es.

 Pretend the table only seats $3$ "people", with $1$ "person" from each planet. Counting clockwise, only the arrangement M, V, E satisfies the given constraints. Therefore, in the actual problem, the members must sit in cycles of M, V, E, but not necessarily with one M, one V, and one E in each cycle(for example, MMVVVE, MVVVEEE, MMMVVVEE all count as cycles). These cycles of MVE must start at seat $1$, since an M is at seat $1$. We simply count the number of arrangements through casework.

 1. The entire arrangement is one cycle- There is only one way to arrange this, MMMMMVVVVVEEEEE

 2. Two cycles - There are 3 Ms, Vs and Es left to distribute among the existing MVEMVE. Using stars and bars, we get $\binom{4}{1}=4$ ways for the members of each planet. Therefore, there are $4^3=64$ ways in total.

 3. Three cycles - 2 Ms, Vs, Es left, so $\binom{4}{2}=6$, making there $6^3=216$ ways total.

 4. Four cycles - 1 M, V, E left, each M can go to any of the four MVE cycles and likewise for V and E, $4^3=64$ ways total

 5. Five cycles - MVEMVEMVEMVEMVE is the only possibility, so there is just $1$ way.

 Combining all these cases, we get $1+1+64+64+216= \boxed{346}$

 





\textbf{Solution 2}

The arrangements must follow the pattern MVEMVE....... where each MVE consists of some Martians followed by some Venusians followed by some Earthlings, for $1, 2, 3, 4,$ or $5$ MVE's. If there are $k$ MVE's, then by stars and bars, there are ${4 \choose k-1}$ choices for the Martians in each block, and the same goes for the Venusians and the Earthlings. Thus, we have $N = 1^3+4^3+6^3+4^3+1^3 = \boxed{346}$ - aops5234

 \textbf{Note} You are probably confused by the solutions. I was too when I first read them. Here is the intuition:Consider the the arrangement of the $M$'s $V$'s and $E$'s: $MMMMM$ | $VVVVV$ | $EEEEE$. Now, we can place a bar in the first Martian's section in $4$ spots, and same with the Venusians and Earthlings.Now for instance, $MM|M|MM$    ||    $V|V|VVV$    ||    $EEE|E|E$ would represent the table ordering MMVEEE MVE MMVVVE, which is a valid ordering. Hence we proceed with the solution and we are done.

 


 


 




	\item (\textit{AIME I 2009}) Consider the set of all triangles $OPQ$ where $O$ is the origin and $P$ and $Q$ are distinct points in the plane with nonnegative integer coordinates $(x,y)$ such that $41x + y = 2009$. Find the number of such distinct triangles whose area is a positive integer.


 
 \textbf{Answer: }600



\textbf{Solution 1 (Matrix's, Determinants)}

Let the two points $P$ and $Q$ be defined with coordinates; $P=(x_1,y_1)$ and $Q=(x_2,y_2)$

 We can calculate the area of the parallelogram with the determinant of the matrix of the coordinates of the two points(shoelace theorem). 

 $\det \left(\begin{array}{c} P \\ Q\end{array}\right)=\det \left(\begin{array}{cc}x_1 &y_1\\x_2&y_2\end{array}\right).$

 Since the triangle has half the area of the parallelogram, we just need the determinant to be even.

 The determinant is

 
 \[(x_1)(y_2)-(x_2)(y_1)=(x_1)(2009-41(x_2))-(x_2)(2009-41(x_1))=2009(x_1)-41(x_1)(x_2)-2009(x_2)+41(x_1)(x_2)=2009((x_1)-(x_2))\]

 Since $2009$ is not even, $x_1-x_2$ must be even, thus the two $x$'s must be of the same parity. Also note that the maximum value for $x$ is $49$ and the minimum is $0$. There are $25$ even and $25$ odd numbers available for use as coordinates and thus there are $(_{25}C_2)+(_{25}C_2)=\boxed{600}$ such triangles.

 





\textbf{Solution 2}

As in Solution 1, let the two points $P$ and $Q$ be defined with coordinates; $P=(x_1,y_1)$ and $Q=(x_2,y_2)$. Let the line $41x+y=2009$ intersect the x-axis at $X$ and the y-axis at $Y$. $X$ has coordinates $(49,0)$, and $Y$ has coordinates $(0,2009)$. As such, there are exactly $50$ lattice points on this line that can be used for $P$ and $Q$.

 


 WLOG, let the x-coordinate of $P$ be less than the x-coordinate of $Q$. Note that $[OPQ]=[OYX]-[OYP]-[OXQ]$. We know that $OY=2009$ and $OX= 49$; as such, $[OYX]=\frac{1}{2} \cdot OY \cdot OX = \frac{1}{2} \cdot 2009 \cdot 49$. In addition, $[OYP]=\frac{1}{2} \cdot 2009 \cdot x_1$ and $[OXQ]=\frac{1}{2} \cdot 49 \cdot y_2$. 

 
Since $2009 \cdot 49$ is odd, the total area of $OYX$ is not an integer; rather, it is of the form $k + \frac{1}{2}$ where $k$ is an integer. To ensure $[OPQ]$ has an integral value, exactly one of $[OPY]$ and $[OQX]$ must have an integral value as well (the other must be of the form $k + \frac{1}{2}$ where $k$ is an integer).

 
Returning to $41x+y=2009$, we notice that integer pairs of $x$ and $y$ that satisfy the equation always have different parities. To satisfy exactly one of $[OPY]$ and $[OQX]$ having an integral area, we must have $x_1$ and $y_2$ having different parities. This is because having an even number for $x_1$ or $y_2$ makes the area of the triangle an integer. We can therefore deduce that $x_1$ and $x_2$ have the same parity. 

 Out of the $50$ usable lattice points for $P$ and $Q$, $25$ have even x-coordinates and $25$ have odd x-coordinates. Since we must pick two points with even x-coordinates or two points with odd x-coordinates, our desired answer is $\binom{25}{2}+\binom{25}{2}=300+300=\fbox{600}$.

 





\textbf{Solution 3 (Analytic geometry)}

As in the solution above, let the two points $P$ and $Q$ be defined with coordinates; $P=(x_1,y_1)$ and $Q=(x_2,y_2)$.

 
If the coordinates of $P$ and $Q$ have nonnegative integer coordinates, $P$ and $Q$ must be lattice points either

 
                \begin{itemize}
                    \item on the nonnegative x-axis
                \end{itemize}


                
                \begin{itemize}
                    \item on the nonnegative y-axis
                \end{itemize}


                
                \begin{itemize}
                    \item in the first quadrant
                \end{itemize}


                We can calculate the y-intercept of the line $41x+y=2009$ to be $(0,2009)$ and the x-intercept to be $(49,0)$.

 Using the point-to-line distance formula, we can calculate the height of $\triangle OPQ$ from vertex $O$ (the origin) to be:

 $\dfrac{|41(0) + 1(0) - 2009|}{\sqrt{41^2 + 1^2}} = \dfrac{2009}{\sqrt{1682}} = \dfrac{2009}{29\sqrt2}$

 Let $b$ be the base of the triangle that is part of the line $41x+y=2009$.

 The area is calculated as: $\dfrac{1}{2}\times b \times \dfrac{2009}{29\sqrt2} = \dfrac{2009}{58\sqrt2}\times b$

 Let the numerical area of the triangle be $k$. 

 So, $k = \dfrac{2009}{58\sqrt2}\times b$

 We know that $k$ is an integer. So, $b = 58\sqrt2 \times z$, where $z$ is also an integer.

 We defined the points $P$ and $Q$ as $P=(x_1,y_1)$ and $Q=(x_2,y_2)$.

 Changing the y-coordinates to be in terms of x, we get:

 $P=(x_1,2009-41x_1)$ and $Q=(x_2,2009-41x_2)$.

 The distance between them equals $b$. 

 Using the distance formula, we get

 $PQ = b = \sqrt{(-41x_2+ 41x_1)^2 + (x_2 - x_1)^2} = 29\sqrt2 \times |x_2 - x_1| = 58\sqrt2\times z$ $(*)$

 WLOG, we can assume that $x_2 > x_1$.

 Taking the last two equalities from the $(*)$ string of equalities and putting in our assumption that $x_2>x_1$, we get

 $29\sqrt2\times (x_2-x_1) = 58\sqrt2\times z$.

 Dividing both sides by $29\sqrt2$, we get

 $x_2-x_1 = 2z$

 As we mentioned, $z$ is an integer, so $x_2-x_1$ is an even integer. Also, $x_2$ and $x_1$ are both positive integers. So, $x_2$ and $x_1$ are between 0 and 49, inclusive. Remember, $x_2>x_1$ as well.

 
                \begin{itemize}
                    \item There are 48 ordered pairs $(x_2,x_1)$ such that their positive difference is 2.
                \end{itemize}


                
                \begin{itemize}
                    \item There are 46 ordered pairs $(x_2,x_1)$ such that their positive difference is 4.
                \end{itemize}


                ...

 
                \begin{itemize}
                    \item Finally, there are 2 ordered pairs $(x_2,x_1)$ such that their positive difference is 48.
                \end{itemize}


                Summing them up, we get that there are $2+4+\dots + 48 = \boxed{600}$ triangles.

 





\textbf{Solution 4}

We present a non-analytic solution; consider the lattice points on the line $41x+y=2009$. The line has intercepts $(0, 2009)$ and $(49, 0)$, so the lattice points for $x=0, 1, \ldots, 49$ divide the line into $49$ equal segments. Call the area of the large triangle $A$. Any triangle formed with the origin having a base of one of these segments has area $A/49$ (call this value $B$) because the height is the same as that of large triangle, and the bases are in the ratio $1:49$. A segment comprised of $n$ small segments (all adjacent to each other) will have area $nB$. Rewriting in terms of the original area, $A=(\frac{1}{2})(49)(2009)$, $B=\frac{2009}{2}$, and $nB=n(\frac{2009}{2})$. It is clear that in order to have a nonnegative integer for $nB$ as desired, $n$ must be even. This is equivalent to finding the number of ways to choose two distinct $x$-values $x_1$ and $x_2$ ($0 \leq x_1, x_2 \leq 49$) such that their positive difference ($n$) is even. Follow one of the previous methods above to choose these pairs and arrive at the answer of 600.

 





\textbf{Solution 5 (Shoelace Theorem)}

Label point $P$ and $Q$ as $(x_1, 2009 - 41x_1)$ and $(x_2, 2009 - 41x_2)$ respectively. Now, we shall use Shoelace Theorem to calculate the area of the triangle. After doing the calculations and simplifications with the coordinates, you will get the area of the triangle to be $\frac{2009|x_1-x_2|}{2}$. Based on this, in order for the area of the triangle to be an integer, $|x_1-x_2|$ needs to be even. The only way for this to happen is that either $x_1$ and $x_2$ are both even are both odd. From the equation from the problem, we can know that $0 \le x_1, x_2 \le 49$. There are $_{25}C_2$ ways to choose 2 even integers in that range or 2 odd numbers in that range. Thus, the answer is $2(_{25}C_2)$ which is $\boxed{600}$.

 





\textbf{Solution 6 (Easier Shoelace)}

Dividing by 41 on both sides yields 
 \[x+ \frac{y}{41} = 49\]

 and because the result is an integer, $\frac{y}{41} = z$ and thus we have
 \[x+z = 49\]and through stars-and-bars the number of ways to count the following integer pairs greater than or equal to 0 is simply just $\binom{50}{1} = 50$. 

 From here, shoelace on the coordinates $(0,0)$, $P(x_1, y_1)$, $Q(x_2,y_2)$ results in $\frac{1}{2} | x_1y_2 - y_1x_2 |$. Also notice that for the resulting area to be an integer we must have that $x_1y_2 \equiv y_1x_2 \pmod{2}$. The only way this is possible is if both $x_1y_2 - y_1x_2$ are either even or odd. There are 25 pairs of odd numbers, and 25 pairs of even numbers out of the 50 possible stars-and-bars cases. Note that choosing the two numbers themselves to give to the coordinates of $P$ and $Q$ respectively make the triangle distinct as order does not matter. Thus the answer is just the number of ways to choose two even coordinates for one of the four distinct values and the remainder is for one of the odd coordinates, and thus our answer is just $2(\binom{25}{2}) = \boxed{600}$.

 




	\item (\textit{AIME I 2009}) In right $\triangle ABC$ with hypotenuse $\overline{AB}$, $AC = 12$, $BC = 35$, and $\overline{CD}$ is the altitude to $\overline{AB}$. Let $\omega$ be the circle having $\overline{CD}$ as a diameter. Let $I$ be a point outside $\triangle ABC$ such that $\overline{AI}$ and $\overline{BI}$ are both tangent to circle $\omega$. The ratio of the perimeter of $\triangle ABI$ to the length $AB$ can be expressed in the form $\frac {m}{n}$, where $m$ and $n$ are relatively prime positive integers. Find $m + n$.


 
 \textbf{Answer: }011



\textbf{Solution 1}

First, note that $AB=37$; let the tangents from $I$ to $\omega$ have length $x$. Then the perimeter of $\triangle ABI$ is equal to 

 




	\item (\textit{AIME I 2009}) The terms of the sequence $\{a_i\}$ defined by $a_{n + 2} = \frac {a_n + 2009} {1 + a_{n + 1}}$ for $n \ge 1$ are positive integers. Find the minimum possible value of $a_1 + a_2$.


 
 \textbf{Answer: }090



\textbf{Solution 1}

Our expression is
 \[a_{n + 2} = \frac {a_n + 2009} {1 + a_{n + 1}}.\]Manipulate this to obtain:
 \[a_{n + 2}a_{n + 1}+a_{n+2}= a_n + 2009.\]Our goal is to cancel terms. If we substitute in $n+1$ for $n,$ we get:
 \[a_{n+3}a_{n+2}+a_{n+3}=a_{n+1}+2009.\]Subtracting these two equations and manipulating the expression yields:
 \[(a_{n+2}+1)(a_{n+3}-a_{n+1})=a_{n+2}-a_n.\]Notice we have the form $a_{k+2}-a_k$ on both sides. Let $b_n=a_{n+2}-a_n.$ Then:
 \[b_{n+1}(a_{n+2}+1)=b_n.\]Notice that since $a_n$ is always an integer, $a_{n+2}+1$ and $b_n$ must also always be an integer. It is also clear that $b_n$ is a multiple of $b_{n+1},$ implying a decreasing sequence.

 However, if the decreasing factor is nonzero, we will eventually have a $b_k$ that is not an integer, contradicting our conditions for $b_n$. Thus, we need either $a_{n+2}+1=0 \Rightarrow a_{n+2}=-1$ (impossible as $a_n$ for all indices must be positive integers) or $b_n=0 \Rightarrow a_{n+2}=a_n.$

 Given this, we want to find the minimum of $a_1+a_2.$ We have, from the problem:
 \[\frac{a_1+2009}{1+a_2}=a_3=a_1 \Rightarrow a_1 a_2 = 2009.\]By AM-GM, to minimize this, we have to make $a_1$ and $a_2$ factors as close as possible. Hence, the smallest possible sum is $41+49=90.$

 





\textbf{Solution 2}

This question is guessable but let's prove our answer

 
 \[a_{n + 2} = \frac {a_n + 2009} {1 + a_{n + 1}}\]

 

 \[a_{n + 2}(1 + a_{n + 1})= a_n + 2009\]

 

 \[a_{n + 2}+a_{n + 2} a_{n + 1}-a_n= 2009\]

 
lets put $n+1$ into $n$ now

 

 \[a_{n + 3}+a_{n + 3} a_{n + 2}-a_{n+1}= 2009\]

 
and set them equal now

 

 \[a_{n + 3}+a_{n + 3} a_{n + 2}-a_{n+1}= a_{n + 2}+a_{n + 2} a_{n + 1}-a_n\]

 

 \[a_{n + 3}-a_{n+1}+a_{n + 3} a_{n + 2}-a_{n + 2} a_{n + 1}= a_{n + 2}-a_n\]

 
let's rewrite it

 

 \[(a_{n + 3}-a_{n+1})(a_{n + 2}+1)= a_{n + 2}-a_n\]

 
Let's make it look nice and let $b_n=a_{n + 2}-a_n$

 

 \[(b_{n+1})(a_{n + 2}+1)= b_n\]

 
Since $b_n$ and $b_{n+1}$ are integers, we can see $b_n$ is divisible by $b_{n+1}$

 
But we can't have an infinite sequence of proper factors, unless $b_n=0$

 
Thus, $a_{n + 2}-a_n=0$

 

 \[a_{n + 2}=a_n\]

 
So now, we know $a_3=a_1$

 

 \[a_{3} = \frac {a_1 + 2009} {1 + a_{2}}\]

 

 \[a_{1} = \frac {a_1 + 2009} {1 + a_{2}}\]

 

 \[a_{1}+a_{1}a_{2} = a_1 + 2009\]

 

 \[a_{1}a_{2} = 2009\]

 
To minimize $a_{1}+a_{2}$, we need $41$ and $49$

 
Thus, our answer $= 41+49=\boxed {090}$

 





\textbf{Solution 3}

If $a_{n} \ne \frac {2009}{a_{n+1}}$, then either
 \[a_{n} = \frac {a_{n}}{1} < \frac {a_{n} + 2009}{1 + a_{n+1}} < \frac {2009}{a_{n+1}}\]

 or

 
 \[\frac {2009}{a_{n+1}} < \frac {2009 + a_{n}}{a_{n+1} + 1} < \frac {a_{n}}{1} = a_{n}\]

 All the integers between $a_{n}$ and $\frac {2009}{a_{n+1}}$ would be included in the sequence. However the sequence is infinite, so eventually there will be a non-integral term.

 So $a_{n} = \frac {2009}{a_{n+1}}$, which $a_{n} \cdot a_{n+1} = 2009$. When $n = 1$, $a_{1} \cdot a_{2} = 2009$. The smallest sum of two factors which have a product of $2009$ is $41 + 49=\boxed {090}$

 


 





\textbf{Solution 4 (BS Solution)}

Essentially you see that it must be an integer for infinite numbers, which doesn't quite seem probable. The most logical explanation is that the sequence repeats, and the numbers in the sequence that repeat are integers. We list out some terms.
 \begin{align*} a_{1} &= a \\ a_{2} &= b \\ a_{3} &=\frac{a+2009}{1+b} \\ a_{4} &=\frac{(b+1)(b+2009)}{a+b+2010} \\ \end{align*}
The terms get more and more wacky, so we just solve for $a,b$ such that $a_{1}=a_{3}$ and $a_{2}=a_{4}.$

 Solving we find both equations end up to the equation $ab=2009$ in which we see to minimize we see that $a = 49$ and $b=41$ or vice versa for an answer of $\boxed{90}.$ This solution is VERY non rigorous and not recommended.

 




	\item (\textit{AIME I 2009}) For $t = 1, 2, 3, 4$, define $S_t = \sum_{i = 1}^{350}a_i^t$, where $a_i \in \{1,2,3,4\}$. If $S_1 = 513$ and $S_4 = 4745$, find the minimum possible value for $S_2$.


 
 \textbf{Answer: }905



\textbf{Solution}

Because the order of the $a$'s doesn't matter, we simply need to find the number of $1$s $2$s $3$s and $4$s that minimize $S_2$. So let $w, x, y,$ and $z$ represent the number of $1$s, $2$s, $3$s, and $4$s respectively. Then we can write three equations based on these variables. Since there are a total of $350$ $a$s, we know that $w + x + y + z = 350$. We also know that $w + 2x + 3y + 4z = 513$ and $w + 16x + 81y + 256z = 4745$. We can now solve these down to two variables:
 \[w = 350 - x - y - z\]Substituting this into the second and third equations, we get
 \[x + 2y + 3z = 163\]and
 \[15x + 80y + 255z = 4395.\]The second of these can be reduced to 
 \[3x + 16y + 51z = 879.\]Now we substitute $x$ from the first new equation into the other new equation.
 \[x = 163 - 2y - 3z\]
 \[3(163 - 2y - 3z) + 16y + 51z = 879\]
 \[489 + 10y + 42z = 879\]
 \[5y + 21z = 195\]Since $y$ and $z$ are integers, the two solutions to this are $(y,z) = (39,0)$ or $(18,5)$.If you plug both these solutions in to $S_2$ it is apparent that the second one returns a smaller value. It turns out that $w = 215$, $x = 112$, $y = 18$, and $z = 5$, so $S_2 = 215 + 4*112 + 9*18 + 16*5 = 215 + 448 + 162 + 80 = \boxed{905}$.

 




	\item (\textit{AIME I 2009}) In triangle $ABC$, $AB = 10$, $BC = 14$, and $CA = 16$. Let $D$ be a point in the interior of $\overline{BC}$. Let points $I_B$ and $I_C$ denote the incenters of triangles $ABD$ and $ACD$, respectively. The circumcircles of triangles $BI_BD$ and $CI_CD$ meet at distinct points $P$ and $D$. The maximum possible area of $\triangle BPC$ can be expressed in the form $a - b\sqrt {c}$, where $a$, $b$, and $c$ are positive integers and $c$ is not divisible by the square of any prime. Find $a + b + c$.


 
 \textbf{Answer: }150



\textbf{Solution 1}

First, by the \href{https://artofproblemsolving.com/wiki/index.php?title=Law\_of\_Cosines}{Law of Cosines}, we have 
 \[\cos BAC = \frac {16^2 + 10^2 - 14^2}{2\cdot 10 \cdot 16} = \frac {256+100-196}{320} = \frac {1}{2},\] so $\angle BAC = 60^\circ$.

 Let $O_1$ and $O_2$ be the circumcenters of triangles $BI_BD$ and $CI_CD$, respectively.  The circles make us think of constant angles, so we first compute 
 \[\angle BO_1D = \angle BO_1I_B + \angle I_BO_1D = 2\angle BDI_B + 2\angle I_BBD.\] Because $\angle BDI_B$ and $\angle I_BBD$ are half of $\angle BDA$ and $\angle ABD$, respectively, the above expression can be simplified to 
 \[\angle BO_1D = \angle BO_1I_B + \angle I_BO_1D = 2\angle BDI_B + 2\angle I_BBD = \angle ABD + \angle BDA.\] Similarly, $\angle CO_2D = \angle ACD + \angle CDA$.  As a result

 
 \begin{align*}\angle CPB &= \angle CPD + \angle BPD \\&= \frac {1}{2} \cdot \angle CO_2D + \frac {1}{2} \cdot \angle BO_1D \\&= \frac {1}{2}(\angle ABD + \angle BDA + \angle ACD + \angle CDA) \\&= \frac {1}{2} (2 \cdot 180^\circ - \angle BAC) \\&= \frac {1}{2} \cdot 300^\circ = 150^\circ.\end{align*}


 Therefore $\angle CPB$ is constant ($150^\circ$). Also, $P$ is $B$ or $C$ when $D$ is $B$ or $C$. Let point $L$ be on the same side of $\overline{BC}$ as $A$ with $LC = LB = BC = 14$; $P$ is on the circle with $L$ as the center and $\overline{LC}$ as the radius, which is $14$. The shortest (and only) distance from $L$ to $\overline{BC}$ is $7\sqrt {3}$.

 When the area of $\triangle BPC$ is the maximum, the distance from $P$ to $\overline{BC}$ has to be the greatest. In this case, it's $14 - 7\sqrt {3}$. The maximum area of $\triangle BPC$ is
 \[\frac {1}{2} \cdot 14 \cdot (14 - 7\sqrt {3}) = 98 - 49 \sqrt {3}\]and the requested answer is $98 + 49 + 3 = \boxed{150}$.

 





\textbf{Solution 2}

From Law of Cosines on $\triangle{ABC}$, 
 \[\cos{A}=\frac{16^2+10^2-14^2}{2\cdot 10\cdot 16}=\frac{1}{2}\implies\angle{A}=60^\circ.\] Setting $\angle{DAC}=\theta$ and doing some angle chasing, we find that $\angle{CI_CD}=\frac{\theta}{2}+90^\circ$ and $\angle{BI_BD}=\frac{\angle{A}-\theta}{2}+90^\circ,$ hence 
 \[\angle{CI_CD}+\angle{BI_BD}=180^\circ+\frac{\angle{A}}{2}=210^\circ.\]Since $CI_CDP$ and $BI_BDP$ are cyclic quadrilaterals, it follows that 
 \[\angle{BPC}=\angle{CPD}+\angle{DPB}=(180^\circ-\angle{CI_CD})+(180^\circ-\angle{BI_BD})=360^\circ-210^\circ=150^\circ.\]Next, applying Law of Cosines on $\triangle{CPB}$,
 \begin{align*} & BC^2=14^2=PC^2+PB^2+2\cdot PB\cdot PC\cdot\frac{\sqrt{3}}{2} \\ & \implies \frac{PC^2+PB^2-196}{PC\cdot PB}=-\sqrt{3} \\ & \implies \frac{PC}{PB}+\frac{PB}{PC}-\frac{196}{PC\cdot PB}=-\sqrt{3} \\ & \implies PC\cdot PB = 196\left(\frac{1}{\frac{PC}{PB}+\frac{PB}{PC}+\sqrt{3}}\right). \end{align*}
 By AM-GM, $\frac{PC}{PB}+\frac{PB}{PC}\geq{2}$, so 
 \[PB\cdot PC\leq 196\left(\frac{1}{2+\sqrt{3}}\right)=196(2-\sqrt{3}).\]Finally, 
 \[[\triangle{BPC}]=\frac12 \cdot PB\cdot PC\cdot\sin{150^\circ}=\frac14 \cdot PB\cdot PC,\]and the maximum area would be $49(2-\sqrt{3})=98-49\sqrt{3},$ so the answer is $\boxed{150}$.

 





\textbf{Solution 3}

Proceed as in Solution 2 until you find $\angle CPB = 150$. The locus of points $P$ that give $\angle CPB = 150$ is a fixed arc from $B$ to $C$ ($P$ will move along this arc as $D$ moves along $BC$) and we want to maximise the area of [$\triangle BPC$]. This means we want $P$ to be farthest distance away from $BC$ as possible, so we put $P$ in the middle of the arc (making $\triangle BPC$ isosceles). We know that $BC=14$ and $\angle CPB = 150$, so $\angle PBC = \angle PCB = 15$. Let $O$ be the foot of the perpendicular from $P$ to line $BC$. Then the area of [$\triangle BPC$] is the same as $7OP$ because base $BC$ has length $14$. We can split $\triangle BPC$ into two $15-75-90$ triangles $BOP$ and $COP$, with $BO=CO=7$ and $OP=7 \tan 15=7(2-\sqrt{3})=14-7\sqrt3$. Then, the area of [$\triangle BPC$] is equal to $7 \cdot OP=98-49\sqrt{3}$, and so the answer is $98+49+3=\boxed{150}$.   

 




\end{enumerate}